% ==========================================
% SECTION 4: PROCESS DISCOVERY AND CONFORMANCE
% ==========================================

This section applies process discovery to construct a formal model from the event log, followed by conformance checking to evaluate both model quality and compliance with legal regulations. A time-based bottleneck analysis and an investigation of batching behavior complete the analysis.

\subsection{Variant Analysis}
\label{sub:control_flow}

A process variant is a unique end-to-end sequence of activities observed in the log. Table~\ref{tab:top_variants} lists the five most frequent variants.

\begin{table}[H]
\centering
\begin{tabular}{lr}
\hline
\textbf{Variant} & \textbf{Cases} \\
\hline
Create Fine $\rightarrow$ Send Fine $\rightarrow$ Insert Fine Notification $\rightarrow$ Add penalty $\rightarrow$ Send for Credit Collection & 56,482 \\
Create Fine $\rightarrow$ Payment & 46,372 \\
Create Fine $\rightarrow$ Send Fine & 20,385 \\
Create Fine $\rightarrow$ Send Fine $\rightarrow$ Insert Fine Notification $\rightarrow$ Add penalty $\rightarrow$ Payment & 9,535 \\
Create Fine $\rightarrow$ Send Fine $\rightarrow$ Insert Fine Notification $\rightarrow$ Add penalty $\rightarrow$ Payment $\rightarrow$ Payment & 3,726 \\
\hline
\end{tabular}
\caption{Top 5 process variants by frequency}
\label{tab:top_variants}
\end{table}

The two dominant variants together account for over 68\% of all completed cases. The most frequent variant ends in credit collection, while the second represents immediate payment without any intermediate steps. This concentration within a small number of variants confirms that the process operates under rigid administrative workflows with limited structural deviation in most cases.

\begin{figure}[H]
    \centering
    \includegraphics[width=0.95\textwidth]{top_variants.png}
    \caption{Frequency distribution of the top 5 process variants.}
    \label{fig:top_variants}
\end{figure}

\subsection{Process Model Discovery}
\label{sub:formal_modeling}

A Petri net was discovered using the Inductive Miner algorithm applied to the ten most frequent process variants. Filtering to the top variants produces a readable model that captures the dominant control-flow logic without being obscured by rare exception paths.

\begin{figure}[H]
    \centering
    \includegraphics[width=0.95\textwidth]{petri_net.png}
    \caption{Petri net discovered via the Inductive Miner on the top 10 variants.}
    \label{fig:petri_net}
\end{figure}

The Petri net reveals the structural logic of the process. After \textit{Create Fine}, the model shows an exclusive choice: a case either proceeds to \textit{Send Fine} (the standard administrative path) or directly to \textit{Payment} (immediate settlement). Along the standard path, a sequence of \textit{Send Fine} $\rightarrow$ \textit{Insert Fine Notification} $\rightarrow$ \textit{Add penalty} leads to a second exclusive choice between \textit{Payment} and \textit{Send for Credit Collection}. The model also captures a loop on the \textit{Payment} activity, reflecting cases where multiple partial payments are recorded. This structure confirms that the process is fundamentally sequential with two key decision points — one at the beginning (immediate payment vs.\ administrative processing) and one at the end (payment vs.\ escalation).

To evaluate how well this model generalizes to the full event log, Token-Based Replay was performed on a random sample of 2,000 cases drawn from the unfiltered log. The replay yielded:

\begin{itemize}
    \item Percentage of fitting traces: 98.70\%
    \item Average trace fitness: 0.9985
\end{itemize}

The high fitness score confirms that the model discovered from the top 10 variants captures nearly all real-world execution paths, including those not explicitly used during discovery.

\subsection{Conformance Checking: Legal and Operational Compliance}
\label{sub:conformance}

Beyond model fitness, five domain-specific compliance rules were defined based on the Italian Codice della Strada (CdS). Each rule was evaluated programmatically against the full event log. The compliance rate is defined as the fraction of applicable cases that do not violate the rule.

\subsubsection{Rule 1: 90-Day Statutory Deadline (Art.\ 201 CdS)}

The Italian Highway Code requires that notification of an offence is sent within 90 days of the violation. The rule checks whether the time between \textit{Create Fine} and \textit{Send Fine} exceeds 90 days.

\begin{itemize}
    \item Applicable cases: 103,987
    \item Violations: 49,772 (average delay in violation cases: 124 days)
    \item \textbf{Compliance rate: 52.14\%}
\end{itemize}

Nearly half of all applicable cases exceed the statutory deadline. This represents a critical compliance failure that may render affected fines legally contestable.

\subsubsection{Rule 2: 60-Day Penalty Grace Period (Art.\ 202\S1, 203\S3 CdS)}

Citizens must be given at least 60 days after notification before a penalty can be added. The rule checks whether \textit{Add penalty} occurs fewer than 60 days after the notification event.

\begin{itemize}
    \item Applicable cases: 79,860
    \item Violations: 0
    \item \textbf{Compliance rate: 100\%}
\end{itemize}

The administration fully respects the 60-day grace period in all recorded cases.

\subsubsection{Rule 3: Due Process — Appeal Prerequisite (Art.\ 203\S1 CdS)}

An appeal to the Prefecture can only be filed if the offender was previously notified. The rule verifies that \textit{Send Appeal to Prefecture} is preceded by either \textit{Insert Fine Notification} or \textit{Send Fine}.

\begin{itemize}
    \item Applicable cases: 4,141
    \item Violations: 7
    \item \textbf{Compliance rate: 99.83\%}
\end{itemize}

Due process is respected in virtually all appeal cases.

\subsubsection{Rule 4: Duty to Inform (Art.\ 204\S2 CdS)}

When the Prefecture issues a decision on an appeal, the administration must notify the offender of the result. The rule checks whether \textit{Receive Result Appeal from Prefecture} is followed by \textit{Notify Result Appeal to Offender}.

\begin{itemize}
    \item Applicable cases: 761
    \item Violations: 43
    \item \textbf{Compliance rate: 94.35\%}
\end{itemize}

In 43 completed cases, the offender was never informed of the Prefecture's decision — a gap that undermines procedural transparency.

\subsubsection{Rule 5: Payment Accuracy (Art.\ 202\S1, 203\S3 CdS)}

Once the statutory debt is settled through payment, no further payment events should occur. The rule detects cases where additional payments are recorded after the cumulative amount has already met or exceeded the required total.

\begin{itemize}
    \item Applicable cases: 69,715
    \item Violations: 7,522
    \item \textbf{Compliance rate: 89.21\%}
\end{itemize}

Over 7,500 cases contain redundant payment events, suggesting either system errors or manual double-processing.

\subsubsection{Compliance Summary}

\begin{table}[H]
    \centering
    \begin{tabular}{llr}
        \toprule
        \textbf{Rule} & \textbf{Legal Basis} & \textbf{Compliance} \\
        \midrule
        90-Day Statutory Deadline & Art.\ 201 CdS & 52.14\% \\
        60-Day Penalty Grace Period & Art.\ 202\S1, 203\S3 CdS & 100.00\% \\
        Due Process (Appeal Prerequisite) & Art.\ 203\S1 CdS & 99.83\% \\
        Duty to Inform & Art.\ 204\S2 CdS & 94.35\% \\
        Payment Accuracy & Art.\ 202\S1, 203\S3 CdS & 89.21\% \\
        \bottomrule
    \end{tabular}
    \caption{Summary of compliance rule evaluation results}
    \label{tab:compliance_summary}
\end{table}

The most critical finding is the widespread violation of the 90-day notification deadline. The remaining rules show high compliance, with the exception of payment accuracy, where redundant payment events suggest a systematic processing issue rather than intentional misconduct.

\subsection{Bottleneck Analysis}
\label{sub:performance_bottlenecks}

Transition durations — the time elapsed between consecutive activities within a case — were computed to identify operational bottlenecks. Figures~\ref{fig:bottlenecks_mean} and~\ref{fig:bottlenecks_median} show the top five transitions ranked by mean and median duration respectively.

\begin{figure}[H]
    \centering
    \includegraphics[width=0.9\textwidth]{bottlenecks_mean.png}
    \caption{Top 5 process bottlenecks by mean transition duration (days).}
    \label{fig:bottlenecks_mean}
\end{figure}

\begin{figure}[H]
    \centering
    \includegraphics[width=0.9\textwidth]{bottlenecks_median.png}
    \caption{Top 5 process bottlenecks by median transition duration (days).}
    \label{fig:bottlenecks_median}
\end{figure}

The mean perspective is dominated by transitions leading to \textit{Send for Credit Collection}, reflecting the long tail of cases that remain unresolved for months or years before escalation. The median perspective highlights a different pattern: the most consistent delay occurs before \textit{Payment}, indicating that the primary bottleneck is driven by citizen response behavior rather than internal administrative processing.

\subsection{Batching Behavior}
\label{sub:batching_analysis}

The Dotted Chart (Figure~\ref{fig:dotted_chart}) reveals a distinctive pattern: events of the same activity type align vertically, meaning dozens of unrelated cases experience the same activity on the same calendar date. This is a clear indicator of batch processing — administrative staff accumulate cases and process them collectively at periodic intervals rather than handling each case immediately upon arrival.

\begin{figure}[H]
    \centering
    \includegraphics[width=0.95\textwidth]{dotted_chart.png}
    \caption{Dotted Chart showing vertical alignment of events as evidence of batching behavior (sample of 500 cases).}
    \label{fig:dotted_chart}
\end{figure}

The Performance Spectrum (Figure~\ref{fig:performance_spectrum}) provides a complementary view. Steep, near-vertical line segments between \textit{Create Fine}, \textit{Send Fine}, and \textit{Insert Fine Notification} confirm that early administrative steps are executed rapidly and in synchronized batches. In contrast, the transitions toward \textit{Add penalty} and \textit{Payment} show shallow, widely dispersed slopes, reflecting long and variable waiting periods.

\begin{figure}[H]
    \centering
    \includegraphics[width=0.95\textwidth]{performance_spectrum.png}
    \caption{Performance Spectrum showing processing velocities across major transitions (150 cases).}
    \label{fig:performance_spectrum}
\end{figure}

\subsection{Organizational Perspective}
\label{sub:organizational}

Process mining typically includes an organizational perspective to analyze resource allocation, handovers of work, and individual workloads. In this dataset, however, the \texttt{org:resource} attribute has a fill rate below one percent and contains only placeholder values. The resource identities were anonymized or omitted during the original data extraction by \textcite{deleoni2015}. As a result, organizational mining — such as constructing social network graphs or computing resource-level performance metrics — cannot be performed on this event log. The analysis therefore focuses on the control-flow, time, and compliance perspectives where the data provides sufficient coverage.
